\documentclass[11pt,a4paper]{report}

% ============================================================
% PACKAGES
% ============================================================
\usepackage[utf8]{inputenc}
\usepackage[T1]{fontenc}
\usepackage{lmodern}
\usepackage[margin=2.5cm]{geometry}
\usepackage{graphicx}
\usepackage{xcolor}
\usepackage{hyperref}
\usepackage{listings}
\usepackage{booktabs}
\usepackage{longtable}
\usepackage{amsmath,amssymb}
\usepackage{enumitem}
\usepackage{fancyhdr}
\usepackage{titlesec}
\usepackage{tocloft}
\usepackage{parskip}

% ============================================================
% STYLING
% ============================================================
\definecolor{codeblue}{RGB}{0,102,204}
\definecolor{codegray}{RGB}{100,100,100}
\definecolor{codebg}{RGB}{248,248,248}
\definecolor{codegreen}{RGB}{0,128,0}
\definecolor{scionblue}{RGB}{25,55,95}
\definecolor{scionaccent}{RGB}{0,120,180}

\hypersetup{
    colorlinks=true,
    linkcolor=scionblue,
    urlcolor=scionaccent,
    citecolor=scionblue,
}

\lstdefinestyle{python}{
    language=Python,
    backgroundcolor=\color{codebg},
    basicstyle=\ttfamily\small,
    keywordstyle=\color{codeblue}\bfseries,
    stringstyle=\color{codegreen},
    commentstyle=\color{codegray}\itshape,
    numbers=left,
    numberstyle=\tiny\color{codegray},
    numbersep=8pt,
    frame=single,
    rulecolor=\color{codegray!30},
    breaklines=true,
    showstringspaces=false,
    tabsize=4,
    xleftmargin=1.5em,
    framexleftmargin=1em,
}

\lstdefinestyle{bash}{
    language=bash,
    backgroundcolor=\color{codebg},
    basicstyle=\ttfamily\small,
    keywordstyle=\color{codeblue},
    frame=single,
    rulecolor=\color{codegray!30},
    breaklines=true,
    showstringspaces=false,
    xleftmargin=1.5em,
    framexleftmargin=1em,
}

\lstset{style=python}

% Header/Footer
\pagestyle{fancy}
\fancyhf{}
\fancyhead[L]{\small\textcolor{codegray}{\leftmark}}
\fancyhead[R]{\small\textcolor{codegray}{Scion Manual}}
\fancyfoot[C]{\thepage}
\renewcommand{\headrulewidth}{0.4pt}

% Chapter title formatting
\titleformat{\chapter}[display]
{\normalfont\huge\bfseries\color{scionblue}}
{\chaptertitlename\ \thechapter}{20pt}{\Huge}

\titleformat{\section}
{\normalfont\Large\bfseries\color{scionblue}}
{\thesection}{1em}{}

\titleformat{\subsection}
{\normalfont\large\bfseries\color{scionblue!80}}
{\thesubsection}{1em}{}

% ============================================================
% DOCUMENT
% ============================================================
\begin{document}

% ------------------------------------------------------------
% TITLE PAGE
% ------------------------------------------------------------
\begin{titlepage}
\centering
\vspace*{3cm}

{\Huge\bfseries\textcolor{scionblue}{Scion}\par}
\vspace{0.3cm}
{\large\textcolor{codegray}{Scientific Computing Infrastructure for Operator Numerics}\par}
\vspace{0.5cm}
{\Large\textcolor{codegray}{A Modular Computational Mechanics Library}\par}

\vspace{2cm}

{\large
Numerical Methods for PDEs $\cdot$ Scientific Machine Learning $\cdot$ Uncertainty Quantification
\par}

\vspace{3cm}

{\large Ioannis Kalogeris\par}
\vspace{0.5cm}
{\normalsize Version 0.1.0\par}
\vspace{0.3cm}
{\normalsize\textcolor{codegray}{February 2026}\par}

\vfill

{\small
Python orchestration with optional Rust/C++ acceleration\\
MIT License
\par}
\end{titlepage}

% ------------------------------------------------------------
% TABLE OF CONTENTS
% ------------------------------------------------------------
\tableofcontents
\newpage

% ============================================================
% CHAPTER 1: INTRODUCTION
% ============================================================
\chapter{Introduction}

\section{Overview}

Scion is a modular computational mechanics library designed for research in numerical methods for partial differential equations (PDEs), scientific machine learning, and uncertainty quantification. The library follows a \emph{Python-first} philosophy: all orchestration happens in Python using NumPy and SciPy, with optional compiled backends in Rust and C++ for performance-critical loops.

\section{Design Principles}

\begin{enumerate}[leftmargin=2em]
    \item \textbf{Python first.} All orchestration in Python. NumPy/SciPy for numerics.
    \item \textbf{Swappable discretisations.} FEM, FDM, FVM, and SPH share the same \texttt{Domain} $\to$ \texttt{LinearSystem} $\to$ \texttt{Solver} pipeline.
    \item \textbf{Library vs.\ projects.} Reusable components live in \texttt{src/scion/}; research workflows live in \texttt{projects/}.
    \item \textbf{Compiled backends optional.} Rust (via maturin/PyO3) and C++ (via pybind11/Eigen) for hot loops. The library works without them.
    \item \textbf{Meshing is external.} Use gmsh or meshio for mesh generation and import.
    \item \textbf{VTK output.} Write \texttt{.vtu} files and visualise in ParaView or interactively with PyVista.
\end{enumerate}

\section{Installation}

Install the core library in editable (development) mode:

\begin{lstlisting}[style=bash]
cd ~/Desktop/ResearchCodes/Scion
pip install -e ".[dev]"
\end{lstlisting}

Install with optional dependency groups as needed:

\begin{lstlisting}[style=bash]
pip install -e ".[jax,dev]"    # JAX + Equinox (for NSF project)
pip install -e ".[ml,dev]"     # PyTorch (for ML subpackage)
pip install -e ".[all]"        # Everything
\end{lstlisting}

% ============================================================
% CHAPTER 2: PROJECT STRUCTURE
% ============================================================
\chapter{Project Structure}

\section{Top-Level Layout}

\begin{lstlisting}[style=bash,numbers=none]
Scion/
    pyproject.toml          # Package config and dependencies
    README.md
    src/scion/         # The library (reusable)
    projects/               # Research projects (consumers)
    tests/                  # Unit tests
    examples/               # Short demos
    benchmarks/             # Performance benchmarks
    docs/                   # Documentation source
\end{lstlisting}

\begin{itemize}[leftmargin=2em]
    \item \textbf{\texttt{src/scion/}} contains all reusable library code. Anything here should be general enough that multiple projects can import it.
    \item \textbf{\texttt{projects/}} contains full research workflows---each project has its own config, training scripts, data pipelines, and checkpoints.
    \item \textbf{\texttt{examples/}} contains short, self-contained scripts demonstrating library usage (e.g.\ \texttt{poisson\_2d.py}).
    \item \textbf{\texttt{tests/}} contains pytest unit tests for the library modules.
\end{itemize}

\section{Library Subpackages}

\begin{lstlisting}[style=bash,numbers=none]
src/scion/
    core/               # Domain, Field, DOF, logging
    mesh/               # Mesh generation and import
    discretization/     # FEM, FDM, FVM, SPH
    materials/          # Constitutive models
    solvers/            # Linear and nonlinear solvers
    boundary/           # Boundary condition handling
    geometry/           # Shape generation, SDF, harmonic maps
    ml/                 # Machine learning
    uq/                 # Uncertainty quantification
    optimization/       # Optimisation methods
    io/                 # VTK/VTU output
    postprocessing/     # PyVista visualisation
    _backends/          # Rust and C++ acceleration
\end{lstlisting}

\section{Projects}

\begin{lstlisting}[style=bash,numbers=none]
projects/
    nsf_domain_mapping/     # Neural Spline Flow: domain -> square
        models/             #   NSF architecture
        config.py           #   Hyperparameters
        dataset.py          #   Data generation pipeline
        Main_nsf.py         #   Full training/eval pipeline
        data/               #   Generated datasets
        checkpoints_nsf/    #   Trained models
\end{lstlisting}

% ============================================================
% CHAPTER 3: CORE SIMULATION
% ============================================================
\chapter{Core Simulation Modules}

\section{\texttt{core} --- Domain, Field, and DOF Management}

The \texttt{core} subpackage provides the foundational abstractions.

\begin{longtable}{lp{9cm}}
\toprule
\textbf{Class/Function} & \textbf{Description} \\
\midrule
\texttt{Domain} & Represents a computational domain with mesh, fields, and boundary conditions \\
\texttt{Field} & Scalar or vector field defined on a mesh \\
\texttt{DOFManager} & Maps mesh nodes to degrees of freedom, handles numbering and constraints \\
\texttt{get\_logger(name)} & Consistent logging across the library \\
\bottomrule
\end{longtable}

\section{\texttt{mesh} --- Mesh Generation and Import}

\begin{longtable}{lp{9cm}}
\toprule
\textbf{Class/Function} & \textbf{Description} \\
\midrule
\texttt{StructuredMesh} & Regular grid mesh generation (Quad4, Hex8) \\
\texttt{UnstructuredMesh} & Triangular/tetrahedral mesh via gmsh \\
\texttt{read\_mesh(path)} & Import mesh files via meshio (.msh, .vtu, .vtk) \\
\bottomrule
\end{longtable}

\section{\texttt{discretization} --- Numerical Schemes}

All discretisation methods inherit a common assembler interface that produces a sparse linear system $\mathbf{K}\mathbf{u} = \mathbf{f}$.

\subsection{FEM --- Finite Element Method}

\begin{longtable}{lp{9cm}}
\toprule
\textbf{Class/Function} & \textbf{Description} \\
\midrule
\texttt{FEMAssembler} & Assembles stiffness matrix and load vector \\
\texttt{Element} & Base class for element types (Tri3, Quad4, Tet4, Hex8) \\
\texttt{GaussQuadrature} & Numerical integration rules \\
\texttt{ShapeFunctions} & Shape function evaluation and gradients \\
\bottomrule
\end{longtable}

\subsection{FDM --- Finite Difference Method}

\begin{longtable}{lp{9cm}}
\toprule
\textbf{Class/Function} & \textbf{Description} \\
\midrule
\texttt{FDMAssembler} & Builds finite difference stencil matrix \\
\texttt{Stencil} & Configurable stencil patterns (3-point, 5-point, 9-point) \\
\bottomrule
\end{longtable}

\section{\texttt{materials} --- Constitutive Models}

\begin{longtable}{lp{9cm}}
\toprule
\textbf{Class} & \textbf{Description} \\
\midrule
\texttt{LinearElastic} & Isotropic linear elasticity (Young's modulus, Poisson's ratio) \\
\texttt{ThermalMaterial} & Thermal conductivity for heat transfer problems \\
\texttt{HyperelasticMaterial} & Base class for nonlinear materials (scaffold) \\
\bottomrule
\end{longtable}

\section{\texttt{solvers} --- Linear and Nonlinear Solvers}

\begin{longtable}{lp{9cm}}
\toprule
\textbf{Class} & \textbf{Description} \\
\midrule
\texttt{DirectSolver} & SciPy sparse direct solver (spsolve) \\
\texttt{ConjugateGradient} & CG iterative solver for SPD systems \\
\texttt{GMRES} & Generalised Minimal Residual for non-symmetric systems \\
\texttt{NewtonRaphson} & Newton--Raphson for nonlinear problems \\
\bottomrule
\end{longtable}

\section{\texttt{boundary} --- Boundary Conditions}

\begin{longtable}{lp{9cm}}
\toprule
\textbf{Class/Function} & \textbf{Description} \\
\midrule
\texttt{DirichletBC} & Prescribed displacement/temperature values \\
\texttt{NeumannBC} & Prescribed traction/flux on boundaries \\
\texttt{apply\_bcs(K, f, bcs)} & Apply boundary conditions to the assembled system \\
\bottomrule
\end{longtable}

\section{\texttt{io} and \texttt{postprocessing}}

\begin{longtable}{lp{9cm}}
\toprule
\textbf{Class/Function} & \textbf{Description} \\
\midrule
\texttt{write\_vtu(path, mesh, fields)} & Write solution to VTU format \\
\texttt{plot\_mesh(mesh)} & Interactive 3D mesh visualisation via PyVista \\
\texttt{plot\_field(mesh, field)} & Scalar/vector field plots with colour maps \\
\bottomrule
\end{longtable}

% ============================================================
% CHAPTER 4: GEOMETRY
% ============================================================
\chapter{Geometry Subpackage}

The \texttt{scion.geometry} subpackage provides tools for generating parameterised 2D shapes, computing signed distance fields, and solving harmonic maps from arbitrary domains to the unit square.

\section{\texttt{geometry.shapes} --- Parameterised Closed Curves}

All shape classes inherit from \texttt{ClosedCurve} and implement \texttt{\_\_call\_\_(t)} for evaluation at parameter $t \in [0, 2\pi]$.

\begin{longtable}{lp{8cm}}
\toprule
\textbf{Class} & \textbf{Description} \\
\midrule
\texttt{ClosedCurve} & Abstract base class: \texttt{\_\_call\_\_(t)}, \texttt{sample(n)}, \texttt{contains(pts)} \\
\texttt{Ellipse} & Axis-aligned or rotated ellipse with semi-axes $a$, $b$ \\
\texttt{Superellipse} & Lam\'e curve: $|x/a|^n + |y/b|^n = 1$ \\
\texttt{FourierBlob} & Radial Fourier perturbation: $r(\theta) = r_0 + \sum a_k \cos(k\theta) + b_k \sin(k\theta)$ \\
\texttt{StarPolygon} & Star polygon with alternating inner/outer radii \\
\texttt{SmoothPolygon} & Polygon with rounded corners via Catmull--Rom splines \\
\texttt{ShapeGenerator} & Factory for generating random shapes with configurable distribution \\
\bottomrule
\end{longtable}

Each class provides a \texttt{random(rng)} class method for random instantiation.

\begin{lstlisting}[caption={Generating and sampling shapes}]
from scion.geometry import ShapeGenerator, Ellipse

# Random shapes with custom distribution
gen = ShapeGenerator(seed=42, shape_distribution={"Ellipse": 0.5, "FourierBlob": 0.5})
shape = gen.generate()

# Specific shape
ellipse = Ellipse(center=(0, 0), a=1.2, b=0.7, theta=0.3)
points = ellipse.sample(500)  # (500, 2) array of boundary points
\end{lstlisting}

\section{\texttt{geometry.sdf} --- Signed Distance Fields}

\begin{longtable}{lp{8cm}}
\toprule
\textbf{Class/Function} & \textbf{Description} \\
\midrule
\texttt{SDFRenderer} & Batch renderer: shape $\to$ normalised SDF image \\
\texttt{render\_shape\_to\_sdf(shape, ...)} & Compute SDF on a regular grid \\
\texttt{sample\_points\_in\_shape(shape, n)} & Rejection sampling of interior points \\
\texttt{normalize\_sdf(sdf, method)} & Normalise SDF to $[-1, 1]$ (tanh, clip, or linear) \\
\bottomrule
\end{longtable}

\section{\texttt{geometry.harmonic\_maps} --- Harmonic Map Solver}

Computes the ground-truth diffeomorphism $\varphi: \Omega \to [-1,1]^2$ by solving two Laplace equations $\Delta u = 0$, $\Delta v = 0$ with boundary conditions mapping $\partial\Omega$ to the square boundary.

\begin{longtable}{lp{8cm}}
\toprule
\textbf{Class/Function} & \textbf{Description} \\
\midrule
\texttt{HarmonicMapSolver} & Solve and cache harmonic maps for a given shape \\
\texttt{compute\_harmonic\_map(shape, ...)} & Full FD solve returning $(u, v)$ grids \\
\texttt{interpolate\_harmonic\_map(...)} & Bilinear interpolation at arbitrary query points \\
\bottomrule
\end{longtable}

% ============================================================
% CHAPTER 5: MACHINE LEARNING
% ============================================================
\chapter{Machine Learning Subpackage}

The \texttt{scion.ml} subpackage provides neural network architectures, training utilities, and surrogate model wrappers built on PyTorch. Install with \texttt{pip install -e ".[ml]"}.

\section{\texttt{ml.networks} --- Neural Network Architectures}

\begin{longtable}{lp{8cm}}
\toprule
\textbf{Class} & \textbf{Description} \\
\midrule
\texttt{MLP} & Fully connected network with configurable depth, width, activation, dropout, and batch normalisation \\
\texttt{MeshGNN} & Graph neural network operating on mesh connectivity. Nodes = mesh nodes, edges = element connections \\
\texttt{UNet} & U-Net encoder--decoder for field-to-field regression on structured grids \\
\bottomrule
\end{longtable}

\begin{lstlisting}[caption={Creating an MLP}]
from scion.ml.networks import MLP

net = MLP(
    layer_sizes=[3, 128, 128, 1],
    activation="gelu",
    dropout=0.1,
)
print(net)  # MLP([3, 128, 128, 1], params=17,153)
\end{lstlisting}

\section{\texttt{ml.training} --- Training Utilities}

\begin{longtable}{lp{8cm}}
\toprule
\textbf{Class} & \textbf{Description} \\
\midrule
\texttt{Trainer} & Generic PyTorch training loop with early stopping, checkpointing, and LR scheduling \\
\texttt{MSELoss} & Standard mean squared error \\
\texttt{RelativeMSELoss} & Relative $L^2$ loss: $\|\hat{y} - y\|^2 / \|y\|^2$ \\
\texttt{PhysicsInformedLoss} & PDE residual loss using automatic differentiation \\
\texttt{WeightedLoss} & Weighted combination of multiple loss terms \\
\bottomrule
\end{longtable}

\section{\texttt{ml.surrogates} --- PDE Surrogate Models}

\begin{longtable}{lp{8cm}}
\toprule
\textbf{Class} & \textbf{Description} \\
\midrule
\texttt{PDESurrogate} & Wraps a trained network as a drop-in replacement for a PDE solver. Includes input/output scaling and batch prediction. \\
\bottomrule
\end{longtable}

% ============================================================
% CHAPTER 6: UNCERTAINTY QUANTIFICATION
% ============================================================
\chapter{Uncertainty Quantification Subpackage}

The \texttt{scion.uq} subpackage provides sampling methods, Bayesian inference tools, polynomial chaos expansions, and Gaussian process emulators.

\section{\texttt{uq.sampling} --- Sampling Methods}

\begin{longtable}{lp{8cm}}
\toprule
\textbf{Function} & \textbf{Description} \\
\midrule
\texttt{monte\_carlo(n, dim, bounds)} & Standard Monte Carlo sampling \\
\texttt{latin\_hypercube(n, dim, bounds)} & Latin Hypercube Sampling for better space-filling \\
\texttt{sobol\_sequence(n, dim, bounds)} & Sobol quasi-random low-discrepancy sequence \\
\texttt{importance\_sampling(...)} & Importance sampling with normalised weights \\
\bottomrule
\end{longtable}

\section{\texttt{uq.bayesian} --- Bayesian Inference}

\begin{longtable}{lp{8cm}}
\toprule
\textbf{Class} & \textbf{Description} \\
\midrule
\texttt{MetropolisHastings} & Random-walk Metropolis--Hastings with configurable proposal \\
\texttt{AdaptiveMCMC} & Adaptive Metropolis with Haario-style covariance adaptation \\
\texttt{RealNVP} & Normalizing flow for variational inference (affine coupling layers) \\
\texttt{MCMCResult} & Container with samples, log-probs, thinning, and burn-in methods \\
\bottomrule
\end{longtable}

\section{\texttt{uq.pce} --- Polynomial Chaos Expansions}

\begin{longtable}{lp{8cm}}
\toprule
\textbf{Class} & \textbf{Description} \\
\midrule
\texttt{PolynomialChaosExpansion} & Non-intrusive PCE via least-squares. Supports Legendre (uniform) and Hermite (Gaussian) bases. Properties: \texttt{mean}, \texttt{variance}, \texttt{sobol\_indices()}. \\
\bottomrule
\end{longtable}

\section{\texttt{uq.gp} --- Gaussian Process Emulation}

\begin{longtable}{lp{8cm}}
\toprule
\textbf{Class} & \textbf{Description} \\
\midrule
\texttt{GaussianProcessEmulator} & Scalar-output GP with RBF/ARD kernel. Returns \texttt{GPPrediction} with mean, std, and variance. \\
\bottomrule
\end{longtable}

% ============================================================
% CHAPTER 7: OPTIMISATION
% ============================================================
\chapter{Optimisation Subpackage}

The \texttt{scion.optimization} subpackage provides gradient-based optimisers, derivative-free methods, Bayesian optimisation, and optimal experimental design.

\section{\texttt{optimization.gradient} --- Gradient-Based Methods}

\begin{longtable}{lp{8cm}}
\toprule
\textbf{Class} & \textbf{Description} \\
\midrule
\texttt{ScipyOptimizer} & Wrapper around \texttt{scipy.optimize.minimize}. Supports L-BFGS-B, SLSQP, trust-constr. Returns \texttt{OptimizationResult} with history. \\
\bottomrule
\end{longtable}

\section{\texttt{optimization.evolutionary} --- Derivative-Free Methods}

\begin{longtable}{lp{8cm}}
\toprule
\textbf{Class} & \textbf{Description} \\
\midrule
\texttt{CMAES} & Covariance Matrix Adaptation Evolution Strategy. Robust for ill-conditioned, non-separable problems up to $\sim$100 parameters. \\
\bottomrule
\end{longtable}

\section{\texttt{optimization.bayesian\_opt} --- Bayesian Optimisation}

\begin{longtable}{lp{8cm}}
\toprule
\textbf{Class/Function} & \textbf{Description} \\
\midrule
\texttt{BayesianOptimizer} & Sequential BO using GP surrogates from \texttt{uq.gp}. Multi-start acquisition optimisation. \\
\texttt{expected\_improvement(...)} & EI acquisition function \\
\texttt{upper\_confidence\_bound(...)} & UCB acquisition function \\
\bottomrule
\end{longtable}

\section{\texttt{optimization.oed} --- Optimal Experimental Design}

\begin{longtable}{lp{8cm}}
\toprule
\textbf{Class} & \textbf{Description} \\
\midrule
\texttt{OptimalExperimentalDesign} & Maximises information criteria (A-, D-, E-optimality) via the Fisher Information Matrix $\mathbf{J}^\top \boldsymbol{\Sigma}^{-1} \mathbf{J}$. \\
\bottomrule
\end{longtable}

% ============================================================
% CHAPTER 8: NSF PROJECT
% ============================================================
\chapter{Project: NSF Domain Mapping}

\section{Objective}

Learn a diffeomorphism $\varphi: \Omega \to [-1,1]^2$ that maps an arbitrary 2D domain to the unit square using Neural Spline Flows (NSF). The map is guaranteed to have positive Jacobian determinant by construction, enabling:
\begin{enumerate}[leftmargin=2em]
    \item Discretise and solve PDEs on the simple square domain
    \item Map the solution back to the original domain via the exact inverse $\varphi^{-1}$
    \item Use the Jacobian determinant for change-of-variables in the PDE
\end{enumerate}

\section{Architecture}

The \texttt{ConditionalNSF} model consists of:
\begin{itemize}[leftmargin=2em]
    \item \textbf{CNN encoder}: Encodes the SDF image of the shape into a latent vector $\mathbf{z} \in \mathbb{R}^{128}$
    \item \textbf{Neural Spline Flow}: 8 coupling layers, each with 16-bin rational quadratic splines, conditioned on $\mathbf{z}$
\end{itemize}

Key properties of the NSF approach:
\begin{itemize}[leftmargin=2em]
    \item Positive Jacobian determinant guaranteed (no regularisation needed)
    \item Exact analytical inverse (not iterative)
    \item Efficient log-determinant computation
\end{itemize}

\section{Pipeline}

The \texttt{Main\_nsf.py} script runs the full pipeline:

\begin{enumerate}[leftmargin=2em]
    \item \textbf{Data generation}: Random shapes $\to$ SDF images + harmonic map ground truth
    \item \textbf{Training}: Supervised learning with interior MSE + boundary MSE + log-det regularisation
    \item \textbf{Evaluation}: MSE, boundary error, Jacobian statistics, inverse consistency
    \item \textbf{Visualisation}: Deformed grids, Jacobian fields, predicted vs.\ ground truth
\end{enumerate}

\begin{lstlisting}[style=bash,caption={Running the NSF project}]
cd projects/nsf_domain_mapping
python Main_nsf.py
\end{lstlisting}

\section{Configuration}

All hyperparameters are set at the top of \texttt{Main\_nsf.py}:

\begin{longtable}{lll}
\toprule
\textbf{Parameter} & \textbf{Default} & \textbf{Description} \\
\midrule
\texttt{n\_layers} & 8 & Number of NSF coupling layers \\
\texttt{n\_bins} & 16 & Spline bins per layer \\
\texttt{hidden\_dim} & 128 & Hidden width in parameter networks \\
\texttt{domain\_bounds} & $(-1.5, 1.5)$ & Input domain bounds \\
\texttt{codomain\_bounds} & $(-1.0, 1.0)$ & Target square bounds \\
\texttt{learning\_rate} & $5 \times 10^{-4}$ & AdamW learning rate \\
\texttt{n\_epochs} & 300 & Training epochs \\
\texttt{batch\_size} & 8 & Samples per gradient step \\
\bottomrule
\end{longtable}

% ============================================================
% CHAPTER 9: BACKENDS
% ============================================================
\chapter{Compiled Backends}

\section{Rust Backend}

The Rust backend uses PyO3 and maturin to expose compiled functions to Python. Build with:

\begin{lstlisting}[style=bash]
cd src/scion/_backends/rust
maturin develop --release
\end{lstlisting}

\section{C++ Backend}

The C++ backend uses pybind11 and Eigen. Build with:

\begin{lstlisting}[style=bash]
cd src/scion/_backends/cpp
mkdir build && cd build
cmake .. && make
\end{lstlisting}

Both backends are optional---the library falls back to pure Python/NumPy/SciPy implementations when they are not available.

% ============================================================
% CHAPTER 10: DEVELOPER GUIDE
% ============================================================
\chapter{Developer Guide}

\section{Running Tests}

\begin{lstlisting}[style=bash]
pytest tests/                    # All tests
pytest tests/test_materials/     # Specific module
pytest -v --tb=short             # Verbose with short tracebacks
\end{lstlisting}

\section{Code Style}

The project uses \texttt{ruff} for linting with a line length of 100 characters:

\begin{lstlisting}[style=bash]
ruff check src/scion/
ruff format src/scion/
\end{lstlisting}

\section{Type Checking}

\begin{lstlisting}[style=bash]
mypy src/scion/
\end{lstlisting}

Pylance in VS Code provides real-time type checking. Recommended setting: ``basic'' mode.

\section{Adding a New Project}

To start a new research project:

\begin{enumerate}[leftmargin=2em]
    \item Create \texttt{projects/my\_project/}
    \item Add a \texttt{launch.json} configuration in \texttt{.vscode/} with the correct \texttt{cwd}
    \item Import from the library: \texttt{from scion.geometry import ShapeGenerator}
    \item Import locally within the project: \texttt{from config import Config}
\end{enumerate}

\section{Git Workflow}

Add these to \texttt{.gitignore}:

\begin{lstlisting}[style=bash,numbers=none]
__pycache__/
*.pyc
*.egg-info/
*.vtu
data/
checkpoints*/
outputs*/
.vscode/
\end{lstlisting}

% ============================================================
% APPENDIX
% ============================================================
\appendix

\chapter{Dependencies}

\begin{longtable}{lll}
\toprule
\textbf{Package} & \textbf{Group} & \textbf{Purpose} \\
\midrule
numpy & core & Array operations \\
scipy & core & Sparse matrices, solvers, sampling \\
meshio & core & Mesh I/O (VTK, gmsh, etc.) \\
pyvista & core & Interactive 3D visualisation \\
vtk & core & VTK data structures \\
matplotlib & core & 2D plotting \\
\midrule
torch & ml & Neural network training (PyTorch) \\
\midrule
jax, jaxlib & jax & Automatic differentiation, JIT compilation \\
equinox & jax & Neural networks in JAX \\
optax & jax & Optimisers for JAX \\
diffrax & jax & ODE/SDE solvers in JAX \\
\midrule
gmsh & meshing & Mesh generation \\
\midrule
pytest & dev & Testing framework \\
ruff & dev & Linting and formatting \\
mypy & dev & Static type checking \\
\bottomrule
\end{longtable}

\end{document}
